\documentclass[full-course-notes.tex]{subfiles}

\begin{document}
\lecturemeta{6}
  {From NFA to a Working Scanner: Subset Construction, DFA Minimization, NFA vs.\ DFA
   Trade-offs, and the Design of a Lexical-Analyzer Generator}
  {CLO 1, CLO 2, CLO 3}
  {\S3.7 (Alg.\,3.20, Subset Construction), \S3.9 (Alg.\,3.36, DFA Minimization); \S3.2
   (Input Buffering), \S3.5, \S3.8 (Design of a Lexical-Analyzer Generator)}
  {[ET] Engineering a Compiler, \S2.4--2.5; [AP] Modern Compiler Implementation, \S2.5}
\lecshort{Subset Construction, Minimization \& Lexer Generators}
\setcounter{section}{0}
\makelecturetitle

\tikzset{
  st/.style={draw, circle, minimum size=0.85cm, font=\small, fill=nuLight, draw=nuNavy},
  fin/.style={st, double, double distance=2pt, fill=dbGreenBg, draw=dbGreen},
  dst/.style={draw, rounded corners, minimum height=0.85cm, minimum width=1.7cm, font=\scriptsize, align=center, fill=nuLight, draw=nuNavy},
  dfin/.style={dst, double, double distance=2pt, fill=dbGreenBg, draw=dbGreen},
  dead/.style={dst, fill=examRedBg, draw=examRed},
  trans/.style={thick},
}

\begin{coreidea}[Lecture Roadmap]
\begin{enumerate}
  \item \textbf{Subset construction} (Algorithm 3.20): the mechanical algorithm that turns any
  NFA into an equivalent DFA, worked all the way through on Lecture~\ref{lec:5}'s 14-state NFA
  for $(a\mid b)^{*}abb$.
  \item \textbf{DFA minimization} (Algorithm 3.36): merging behaviorally-identical states to
  produce the smallest possible DFA for the same language -- applied to the very DFA we just
  built.
  \item \textbf{NFA vs.\ DFA, revisited}: same recognizing power (Lecture~\ref{lec:5} already
  proved this constructively), but a real space/time trade-off in practice -- worth
  understanding concretely, not just as a slogan.
  \item Zooming out from theory to engineering: \textbf{how a lexical-analyzer generator is
  actually built} -- combining many token patterns into one automaton, implementing maximal
  munch by simulating a DFA, resolving pattern ties by rule order, and buffering the input
  stream efficiently. This is the generic architecture that \emph{any} scanner-generator tool
  (Flex included) implements under the hood.
\end{enumerate}
This lecture completes the pipeline Lecture~\ref{lec:3}'s big picture
(\Sref{sec:l3-big-picture}) promised: regular expression $\to$ NFA (done in
Lecture~\ref{lec:5}) $\to$ DFA $\to$ minimized DFA $\to$ transition table $\to$ an actual,
running lexer (all done here). By the end of today you will have followed one regular
expression all the way from formal notation to working scanner code, and you will understand
exactly what a tool like Flex is doing on your behalf when you hand it a \texttt{.l} file --
the concrete Flex syntax itself is covered separately in the supplementary Flex notes, not
here.
\end{coreidea}

\section{From an NFA to a DFA: The Subset Construction}
\label{sec:l6-subset-construction}

\dbtag Lecture~\ref{lec:5} left us with a working NFA, but simulating an NFA directly is
awkward: at every step there can be several possible next states, so a naive simulator would
have to explore every possible sequence of choices to know whether \emph{any} of them accepts
(\Sref{sec:l5-dfa-vs-nfa}'s exam-focus box already flagged this). The \textbf{subset
construction} sidesteps the problem entirely with one elegant idea: instead of asking ``which
\emph{one} state am I in?'', track the entire \emph{set} of NFA states that are reachable so
far, all at once. A set of NFA states is itself just one state of a new automaton -- and moving
between these sets, one input symbol at a time, turns out to be completely deterministic.

\subsection{Two Helper Operations}
\label{sec:l6-subset-helpers}

\begin{defbox}{$\varepsilon$-closure and \texttt{move}}
Let $T$ be a set of NFA states.
\begin{itemize}
  \item $\varepsilon\text{-closure}(T)$ is the set of all NFA states reachable from \emph{any}
  state in $T$ using \textbf{zero or more} $\varepsilon$-transitions alone -- ``every state you
  could drift into for free, without reading any input.'' Every state in $T$ is trivially in
  its own $\varepsilon$-closure (zero transitions is allowed).
  \item $\text{move}(T, a)$ is the set of NFA states reachable by taking \textbf{exactly one}
  transition labeled $a \in \Sigma$ from \emph{some} state in $T$ -- ``every state you could be
  in immediately after consuming one $a$, before taking any further $\varepsilon$-transitions.''
\end{itemize}
Every DFA state the subset construction builds is, underneath, one of these
$\varepsilon\text{-closure}$ sets -- which is exactly why DFA states are conventionally drawn or
labeled as \emph{sets} of NFA state numbers throughout this section.
\end{defbox}

\begin{coreidea}[Algorithm 3.20 -- Subset Construction]
\begin{enumerate}
  \item \textbf{Start state.} Compute $D_0 = \varepsilon\text{-closure}(\{n_0\})$, where $n_0$
  is the NFA's start state. $D_0$ becomes the DFA's start state, and is added to a worklist of
  DFA states still needing their outgoing transitions computed.
  \item \textbf{Process the worklist.} While the worklist is nonempty, remove one unmarked DFA
  state $T$ from it. For \emph{every} input symbol $a \in \Sigma$:
  \begin{enumerate}[label=(\alph*)]
    \item Compute $U = \varepsilon\text{-closure}(\text{move}(T, a))$.
    \item If $U$ is empty, there is no transition on $a$ from $T$ -- skip it (this becomes an
    implicit ``error''/dead transition, \Sref{sec:l6-dead-states}).
    \item If $U$ is a \emph{new} set not seen before, add it as a new DFA state and place it on
    the worklist.
    \item Either way, record the transition $\delta_{\mathrm{DFA}}(T, a) = U$.
  \end{enumerate}
  \item \textbf{Accepting states.} A DFA state $T$ (itself a set of NFA states) is accepting if
  and only if $T$ contains \emph{at least one} accepting NFA state -- one successful NFA path
  reaching an accept state is exactly what ``the set contains an accepting state'' captures.
  \item \textbf{Termination.} Since the NFA has only finitely many states, it has only finitely
  many \emph{subsets} of states ($2^{|Q|}$ of them, at most) -- so this process, which only ever
  creates DFA states out of such subsets, is guaranteed to terminate; it cannot generate new DFA
  states forever.
\end{enumerate}
\end{coreidea}

\begin{examfocus}
The single most common hand-execution mistake: forgetting to take the $\varepsilon$-closure
\emph{after} \texttt{move}, not just at the very start. Every new DFA state is
$\varepsilon\text{-closure}(\text{move}(T,a))$, \emph{not} just $\text{move}(T,a)$ -- if you skip
the closure step, you will silently drop reachable NFA states and build an incorrect DFA that
rejects strings it should accept.
\end{examfocus}

\subsection{Worked Example: Building the DFA for $(a\mid b)^{*}abb$}
\label{sec:l6-subset-worked}

\dbtag We continue directly from Lecture~\ref{lec:5}'s 14-state NFA
(\Sref{sec:l5-thompson-worked}): start state 7, single accept state 14, and the $\varepsilon$-
and symbol-transitions built stage by stage in that lecture. Applying subset construction:

\begin{center}
\renewcommand{\arraystretch}{1.3}
\small
\begin{tabularx}{\textwidth}{@{}l X l l@{}}
\toprule
\textbf{DFA State} & \textbf{Set of NFA States (after $\varepsilon$-closure)} & \textbf{On $a$}
& \textbf{On $b$} \\
\midrule
$A$ (start) & $\{7, 5, 3, 1, 8, 9\}$ -- everything reachable from 7 with no input at all & $B$ &
$C$ \\
$B$ & $\varepsilon\text{-closure}(\text{move}(A,a)) = \{2, 10, 6, 11, 8, 5, 1, 3, 9\}$ & $B$ & $D$ \\
$C$ & $\varepsilon\text{-closure}(\text{move}(A,b)) = \{4, 6, 8, 9, 5, 3, 1\}$ & $B$ & $C$ \\
$D$ & $\varepsilon\text{-closure}(\text{move}(B,b)) = \{4, 12, 6, 13, 8, 5, 1, 3, 9\}$ & $B$ & $E$ \\
$E$ (accepting) & $\varepsilon\text{-closure}(\text{move}(D,b)) = \{4, 14, 6, 8, 5, 1, 3, 9\}$ & $B$
& $C$ \\
\bottomrule
\end{tabularx}
\end{center}

\begin{center}
\resizebox{\textwidth}{!}{%
\begin{tikzpicture}[node distance=3.1cm, >=Latex]
  \node[dfin] (A) {$A$\\$\{7,5,3,1,8,9\}$};
  \node[dst, right=of A] (B) {$B$\\$\{2,10,6,11,8,5,1,3,9\}$};
  \node[dst, below=2.1cm of A] (C) {$C$\\$\{4,6,8,9,5,3,1\}$};
  \node[dst, right=of B] (D) {$D$\\$\{4,12,6,13,8,5,1,3,9\}$};
  \node[dfin, right=of D] (E) {$E$\\$\{4,14,6,8,5,1,3,9\}$};
  \draw[trans, ->] (-1.6,0) -- (A);
  \draw[trans, ->] (A) -- node[above, font=\scriptsize]{$a$} (B);
  \draw[trans, ->] (A) -- node[left, font=\scriptsize]{$b$} (C);
  \draw[trans, ->] (B) to[out=60, in=120, looseness=5] node[above, font=\scriptsize]{$a$} (B);
  \draw[trans, ->] (B) -- node[above, font=\scriptsize]{$b$} (D);
  \draw[trans, ->] (D) -- node[above, font=\scriptsize]{$b$} (E);
  \draw[trans, ->] (D) to[bend right=35] node[below, font=\scriptsize]{$a$} (B);
  \draw[trans, ->] (E) to[bend left=45] node[above, font=\scriptsize]{$a$} (B);
  \draw[trans, ->] (E) to[out=200, in=340, looseness=1.1] node[below, font=\scriptsize, pos=0.5]{$b$} (C);
  \draw[trans, ->] (C) to[bend right=25] node[below, font=\scriptsize]{$a$} (B);
  \draw[trans, ->] (C) to[out=250, in=290, looseness=5] node[below, font=\scriptsize]{$b$} (C);
\end{tikzpicture}}
\end{center}

\begin{examplebox}{Reading This DFA in Plain English}
Only state $E$ is accepting, and $E$ is reached only by seeing an $a$ then two $b$'s in a row
($D \xrightarrow{b} E$, and $D$ is reached only via $B\xrightarrow{b}D$, and $B$ is reached only
by having just read an $a$) -- in other words, exactly ``the last three characters read were
\texttt{abb}.'' This is a direct, mechanical confirmation that the DFA really does recognize
$(a\mid b)^{*}abb$: no hand-waving, just five states and the two helper operations above, run to
completion.
\end{examplebox}

\begin{examfocus}
Notice the NFA had \textbf{14 states} but the DFA has only \textbf{5}. This is \emph{not} a
general rule -- subset construction can just as easily \emph{blow up} the state count
(\Sref{sec:l6-nfa-vs-dfa-tradeoffs} works through why) -- it simply happened to shrink here
because many of this NFA's 14 states were $\varepsilon$-only ``glue'' states that collapse into
the same reachable set. Do not assume the DFA is always smaller; always run the algorithm.
\end{examfocus}

\subsection{Dead States and the Implicit Error Transition}
\label{sec:l6-dead-states}

\dbtag Step 2(b) of the algorithm above quietly skips any $\varepsilon\text{-closure}(\text{move}(T,a))$
that comes out empty. In a real scanner, ``no valid transition on this symbol'' cannot simply be
ignored -- it is exactly how a lexical error, or the end of the current token, gets detected. The
standard way to make this explicit is to add one extra, non-accepting \textbf{dead state} $D_{\text{err}}$
with a self-loop on every symbol, and route every missing transition to it.

\begin{center}
\begin{tikzpicture}[node distance=3.8cm, >=Latex]
  \node[dfin] (A) {$A$ (accepting)};
  \node[dead, right=of A] (dead) {$D_{\text{err}}$};
  \draw[trans, ->] (-1.6,0) -- (A);
  \draw[trans, ->] (A) -- node[above, font=\scriptsize, align=center]{\emph{any symbol with no other transition}} (dead);
  \draw[trans, ->] (dead) to[out=45, in=-45, looseness=6] node[right, font=\scriptsize]{\emph{any symbol}} (dead);
\end{tikzpicture}
\end{center}

Reaching $D_{\text{err}}$ means the automaton is certain no continuation can ever lead back to
an accepting state -- so a real scanner treats arrival at $D_{\text{err}}$ as the signal to stop
extending the current match immediately, without even reading further input (there is no reason
to; every path from $D_{\text{err}}$ dead-ends). \Sref{sec:l6-lexgen-maximal-munch} shows exactly
how this plugs into maximal-munch matching.

\section{DFA Minimization}
\label{sec:l6-minimization}

\dbtag Subset construction guarantees a DFA, but not the \emph{smallest} DFA for the language --
different DFA states can behave identically in every possible future, making some of them
redundant. \textbf{DFA minimization} (DB Algorithm 3.36) finds and merges exactly these
redundant states, producing the unique smallest DFA (up to renaming states) recognizing the same
language.

\begin{defbox}{Distinguishable States}
Two DFA states $p$ and $q$ are \textbf{distinguishable} if there exists \emph{some} input
string $w$ such that following $w$ from $p$ ends in an accepting state while following $w$ from
$q$ does not (or vice versa) -- i.e.\ some future input reveals a difference in behavior between
them. States that are \emph{not} distinguishable by \emph{any} string are
\textbf{indistinguishable}, and can be safely merged into a single state without changing which
strings the automaton accepts.
\end{defbox}

\begin{coreidea}[Algorithm 3.36 -- Partition Refinement]
Minimization works by repeatedly \emph{refining} a partition of the states into groups that are
\emph{provisionally} treated as indistinguishable, splitting a group apart the moment a symbol
is found that sends its members to genuinely different groups.
\begin{enumerate}
  \item \textbf{Initial partition.} Split all states into exactly two groups: the accepting
  states $F$, and the non-accepting states $Q \setminus F$ (including the dead state, if one was
  added). This is the coarsest split that could possibly be correct, since an accepting and a
  non-accepting state are always distinguishable by the empty string $w = \varepsilon$ alone.
  \item \textbf{Refine, repeatedly.} For each current group $G$ and each input symbol $a$, check
  where each state in $G$ goes on $a$. If every state in $G$ lands in the \emph{same} other
  group on every symbol, $G$ stays together. If some states in $G$ land in a \emph{different}
  group than others (on some symbol $a$), split $G$ into sub-groups accordingly -- states now
  provably distinguishable (via that one symbol $a$) must never end up merged.
  \item \textbf{Repeat until stable.} Keep refining until one full pass over every group and
  every symbol produces \emph{no} further splits -- at that point the partition has stabilized:
  every two states remaining in the same group really are indistinguishable by \emph{any}
  string, not just the symbols checked so far.
  \item \textbf{Build the minimized DFA.} Each final group becomes one state of the minimized
  DFA. The start state is the group containing the original start state; a group is accepting
  if it contains (any, and therefore every) original accepting state; transitions are inherited
  directly, since every state in a stable group agrees on where each symbol leads.
  \item \textbf{Drop the dead state.} If minimization merges everything unreachable/hopeless
  into the dead-state group, that group is simply omitted from the final diagram -- a missing
  transition in the minimized DFA is understood to mean ``go to the (implicit) dead state and
  reject,'' exactly as \Sref{sec:l6-dead-states} already described.
\end{enumerate}
\end{coreidea}

\begin{examfocus}
The refinement in step 2 must be applied \emph{simultaneously} across every symbol before
deciding a group is stable -- a group only survives a round if it agrees on \emph{every} symbol
in $\Sigma$, not just one. A common hand-execution error is stopping as soon as one symbol
doesn't split a group, without checking the rest of $\Sigma$ too.
\end{examfocus}

\subsection{Worked Example: Minimizing the $(a\mid b)^{*}abb$ DFA}
\label{sec:l6-minimization-worked}

\dbtag Apply the algorithm to the 5-state DFA built in \Sref{sec:l6-subset-worked}
($A,B,C,D,E$; only $E$ accepting):

\begin{center}
\renewcommand{\arraystretch}{1.3}
\small
\begin{tabularx}{\textwidth}{@{}l X@{}}
\toprule
\textbf{Round} & \textbf{Partition} \\
\midrule
Initial & $G_1 = \{A,B,C,D\}$ (non-accepting), $G_2 = \{E\}$ (accepting) \\
Round 1 & Check $G_1$'s members on each symbol. On $a$: $A\to B$, $B\to B$, $C\to B$, $D\to B$
-- all land in $G_1$ (specifically all in $B$'s group so far, fine). On $b$: $A\to C \in G_1$,
$B\to D \in G_1$, $C\to C \in G_1$, but $D\to E \in G_2$ -- \textbf{different group!} $D$ must
split off: $G_1$ becomes $G_{1a}=\{A,B,C\}$ and $G_{1b}=\{D\}$. \\
Round 2 & Re-check $G_{1a}=\{A,B,C\}$: on $a$, all three go to $B \in G_{1a}$ -- consistent. On
$b$: $A\to C\in G_{1a}$, $B\to D\in G_{1b}$, $C\to C\in G_{1a}$ -- \textbf{$B$ disagrees!} Split
again: $G_{1a}$ becomes $\{A,C\}$ and $\{B\}$. \\
Round 3 & Re-check $\{A,C\}$: on $a$, both go to $B$'s group -- consistent. On $b$: $A\to C$,
$C\to C$ -- both land in the \emph{same} group $\{A,C\}$ itself -- consistent. No further split
needed for this group. All other groups are singletons, so they are trivially stable. \\
Stable & $\{A,C\}$, $\{B\}$, $\{D\}$, $\{E\}$ -- four groups, no further splits possible. \\
\bottomrule
\end{tabularx}
\end{center}

\begin{examplebox}{Why $A$ and $C$ Merge}
Intuitively: $A$ (``seen nothing useful yet, or just finished a full \texttt{abb} and are
starting over'') and $C$ (``just saw a \texttt{b} that isn't continuing an \texttt{ab}
prefix'') behave \emph{identically} from here on -- from both, an $a$ starts a fresh attempt at
matching \texttt{abb}, and a $b$ leaves you in exactly the same ``waiting'' situation. No future
string can ever tell $A$ and $C$ apart, which is precisely the formal definition of
indistinguishable states being applied.
\end{examplebox}

\begin{center}
\resizebox{\textwidth}{!}{%
\begin{tikzpicture}[node distance=3.3cm, >=Latex]
  \node[dfin] (AC) {$\{A,C\}$};
  \node[dst, right=of AC] (B) {$\{B\}$};
  \node[dst, right=of B] (D) {$\{D\}$};
  \node[dfin, right=of D] (E) {$\{E\}$};
  \draw[trans, ->] (-1.4,0) -- (AC);
  \draw[trans, ->] (AC) to[out=135, in=225, looseness=6] node[left, font=\scriptsize]{$b$} (AC);
  \draw[trans, ->] (AC) -- node[above, font=\scriptsize]{$a$} (B);
  \draw[trans, ->] (B) to[out=60, in=120, looseness=5] node[above, font=\scriptsize]{$a$} (B);
  \draw[trans, ->] (B) -- node[above, font=\scriptsize]{$b$} (D);
  \draw[trans, ->] (D) -- node[above, font=\scriptsize]{$b$} (E);
  \draw[trans, ->] (D) to[bend right=35] node[below, font=\scriptsize]{$a$} (B);
  \draw[trans, ->] (E) to[bend left=45] node[above, font=\scriptsize]{$a$} (B);
  \draw[trans, ->] (E) to[out=200, in=340, looseness=1.1] node[below, font=\scriptsize, pos=0.5]{$b$} (AC);
\end{tikzpicture}}
\end{center}

\begin{center}
\small\textit{The minimized DFA: 4 states instead of 5 -- $A$ and $C$ have collapsed into one
state, and every transition is inherited unchanged (e.g.\ $E\xrightarrow{b}C$ becomes
$E\xrightarrow{b}\{A,C\}$). This is provably the \emph{smallest possible} DFA for
$(a\mid b)^{*}abb$: DB Algorithm 3.36 guarantees minimality, not just a smaller result.}
\end{center}

\begin{examfocus}
``Smallest possible'' is a real, provable guarantee, not just ``smaller than before.'' Two
different minimized DFAs for the \emph{same} language are always identical up to renaming
states -- there is exactly one minimal DFA (isomorphism aside), which is why minimization is
sometimes described as computing a \emph{canonical form} for a regular language.
\end{examfocus}

\section{NFA vs.\ DFA, Revisited: Power and Performance}
\label{sec:l6-nfa-vs-dfa-tradeoffs}

\dbtag Lecture~\ref{lec:5} and \Sref{sec:l3-chomsky-hierarchy} already established that NFAs
and DFAs recognize \emph{exactly} the same class of languages -- the subset construction above
is the constructive proof of that fact, run by hand. Equal recognizing power does \emph{not}
mean equal practical cost, though, and this distinction is worth making concrete now that you
have built both an NFA and a DFA for the same language.

\begin{center}
\renewcommand{\arraystretch}{1.4}
\begin{tabularx}{\textwidth}{@{}l X X@{}}
\toprule
 & \textbf{NFA (simulated directly)} & \textbf{DFA (pre-built, table-driven)} \\
\midrule
\textbf{States tracked per input symbol} & A \emph{set} of up to $|Q_{\text{NFA}}|$ states must
be tracked and updated at every step (recomputing an $\varepsilon$-closure each time) &
Exactly \textbf{one} state at all times -- a single array index \\
\textbf{Work per input symbol} & $O(|Q_{\text{NFA}}|)$ or worse, since every state in the
current set may need its own transitions and closure explored & $O(1)$ -- one table lookup,
\texttt{table[state][symbol]} \\
\textbf{Number of automaton states} & Exactly as many as Thompson's construction produced --
linear in the size of the regular expression & Can be \textbf{exponentially larger}
($2^{|Q_{\text{NFA}}|}$ in the worst case) -- see the classic blow-up example below \\
\textbf{When the cost is paid} & Every single run, on every input, pays the per-symbol
simulation overhead & Paid \textbf{once}, offline, when the DFA (and its table) is built; every
subsequent run is just table lookups \\
\bottomrule
\end{tabularx}
\end{center}

\begin{examplebox}{The Classic Exponential Blow-up}
Consider the regular expression $(a\mid b)^{*}a(a\mid b)^{n}$ for some fixed $n$ -- ``the
$n$-th-from-last character is an \texttt{a}.'' Its NFA (via Thompson's construction) has only
$O(n)$ states. But its \emph{minimal} DFA needs $2^{n}$ states: intuitively, the DFA must
remember the \emph{last $n$ characters seen} (to know, once $n$ more characters arrive, whether
that remembered character really was $n$-from-last), and there are $2^{n}$ possible values for
``the last $n$ characters,'' each requiring a genuinely distinguishable state. This is a real,
provable worst case, not a contrived edge case -- it is exactly why some lexer generators offer
a choice between DFA-based and NFA-based matching for pathological patterns.
\end{examplebox}

\begin{coreidea}[Why Lexer Generators Overwhelmingly Choose the DFA Side]
For \emph{ordinary} lexical tokens (identifiers, numbers, keywords, operators), the exponential
worst case above essentially never triggers in practice -- real token patterns are small and
well-behaved, so their DFAs stay modest in size. Given that, the DFA's $O(1)$-per-character,
pre-paid-once cost is a clear win for a component like a lexer that runs on \emph{every single
character} of every source file, potentially billions of times across a large codebase. This is
precisely why every mainstream lexer generator (Flex included) compiles down to a DFA rather
than simulating an NFA at scan time, even though building that DFA (subset construction, then
minimization) is itself extra one-time work at generator-build time.
\end{coreidea}

\begin{examfocus}
Contrast this with \Sref{sec:l4-regex-in-practice}'s discussion of backtracking regex engines
(Perl, Python, JavaScript): those engines simulate something closer to an NFA \emph{with
backreferences} at match time, which is exactly why they can suffer ReDoS
(\Sref{sec:l4-regex-in-practice}) -- unlike a lexer generator, they cannot always fall back to a
pre-built DFA, because backreferences push the language being matched outside the regular
languages entirely (\Sref{sec:l4-regex-in-practice} again). A lexer's patterns, by contrast, are
always \emph{genuinely} regular by construction, so the DFA route is always available.
\end{examfocus}

\begin{examfocus}
Lazy (``on-the-fly'') DFA construction is the practical middle ground worth knowing by name:
tools like Google's RE2 and many \texttt{grep} implementations build DFA states
\emph{on demand}, the first time they are actually reached during a real match, and cache them
for reuse -- getting the DFA's fast steady-state lookup without paying to build every
theoretically-possible state up front. This matters most for patterns where the full DFA would
be large but any single input only ever visits a small fraction of its states.
\end{examfocus}

\section{Designing a Lexical-Analyzer Generator}
\label{sec:l6-lexgen-design}

\dbtag Everything so far in this lecture -- and the two before it -- has been about turning
\emph{one} regular expression into \emph{one} DFA. A real lexer, though, must recognize
\emph{dozens} of different token patterns (\texttt{ID}, \texttt{NUM}, every keyword, every
operator, \dots) at once, decide which pattern produced the \emph{longest} match at every
position (maximal munch, \Sref{sec:l3-beyond-db}), and hand the winning token off to the parser.
This section builds that generic architecture -- exactly what a lexer-generator tool constructs
for you, and exactly what \emph{any} such tool (Flex included) is doing under the hood, whatever
its concrete input syntax looks like. \emph{(The concrete \texttt{.l}-file syntax for Flex
specifically is covered in the separate supplementary Flex notes, not here.)}

\subsection{Combining Many Patterns Into One Automaton}
\label{sec:l6-lexgen-combining}

\dbtag A lexer-generator's input is fundamentally a list of \textbf{(pattern, action)} pairs --
one regular expression per token, paired with the code that should run when that token is
matched (typically: build the right \texttt{<token-name, attribute-value>} pair and return it).
The combining step is a direct generalization of the union rule from Thompson's construction
(\Sref{sec:l5-thompson-rules}):

\begin{enumerate}
  \item Build (or imagine) a separate small NFA fragment $N(r_i)$ for each individual pattern
  $r_1, r_2, \dots, r_k$, exactly as in Lecture~\ref{lec:5}.
  \item Introduce one new overall start state, with an $\varepsilon$-transition to \emph{every}
  fragment's own start state -- precisely the union construction, generalized from two
  fragments to $k$.
  \item \textbf{Critically, do \emph{not} merge the fragments' accept states into one shared
  accept state} the way an ordinary two-way union would. Instead, each fragment $N(r_i)$ keeps
  its \emph{own} accept state, tagged with \emph{which} token $r_i$ corresponds to.
  \item Run subset construction (\Sref{sec:l6-subset-construction}) on this combined NFA exactly
  as before. A resulting DFA state is accepting if its underlying set of NFA states contains
  \emph{any} tagged accept state -- and it is tagged with \emph{that} token (or, if more than
  one tagged accept state appears in the same set, with whichever pattern is listed
  \emph{first} in the original specification -- see \Sref{sec:l6-lexgen-ties} below).
\end{enumerate}

\begin{examplebox}{A Miniature Combined Lexer: Two Patterns}
Suppose a lexer needs only two token patterns: \texttt{KEYWORD\_IF} for the literal string
\texttt{if}, and \texttt{ID} for \texttt{letter(letter|digit)*}. Thompson's construction gives
two small fragments; the combining step wires a shared start state into both, with each
fragment's own accept state tagged \texttt{IF} or \texttt{ID} respectively. Subset construction
then merges these into one DFA. Scanning the input \texttt{if5}: after \texttt{i}, then
\texttt{f}, the DFA is in a state whose underlying NFA-state set includes \emph{both} the
\texttt{IF} fragment's accept state \emph{and} the \texttt{ID} fragment's accept state
simultaneously (since \texttt{if} is a valid lexeme for \emph{both} patterns) -- but reading one
more character, \texttt{5}, only the \texttt{ID} fragment has anywhere left to go, so the DFA
continues and eventually accepts \texttt{if5} as a single \texttt{ID} token, illustrating
exactly why maximal munch (next subsection) has to keep reading past an early, shorter
match.
\end{examplebox}

\subsection{Maximal Munch, Implemented as DFA Simulation}
\label{sec:l6-lexgen-maximal-munch}

\dbtag \Sref{sec:l3-beyond-db} introduced the longest-match rule informally. With a combined DFA
in hand, it has a precise, mechanical implementation:

\begin{coreidea}[The Maximal-Munch Scanning Loop]
\begin{enumerate}
  \item Starting from the DFA's start state and the current input position, repeatedly follow
  transitions on successive input characters, \textbf{remembering the most recent point at
  which the current state was accepting} (and which token it was tagged with).
  \item Keep advancing as long as a transition exists for the next character -- \emph{even past}
  an accepting state, since a longer match might still be possible (exactly the \texttt{if5}
  example above: the DFA is accepting for \texttt{IF} right after \texttt{if}, but a real
  scanner does not stop there).
  \item Stop advancing the moment either (a) the input is exhausted, or (b) the next character
  has no valid transition from the current state -- i.e.\ the only remaining move is to the
  implicit dead state $D_{\text{err}}$ from \Sref{sec:l6-dead-states}. There is never a reason
  to enter $D_{\text{err}}$ itself: once every remaining transition would lead there, the scan
  is over.
  \item \textbf{Roll back} to the \emph{last remembered accepting state} (step 1), emit a token
  for whichever pattern that state was tagged with, and resume scanning from the character right
  after that point.
  \item If the current state was \emph{never} accepting at any point during the scan (no
  remembered position exists at all), this is a genuine lexical error
  (\Sref{sec:l3-lexical-errors}): report it and invoke panic-mode recovery
  (\Sref{sec:l3-error-recovery}).
\end{enumerate}
\end{coreidea}

\begin{examplebox}{Tracing \texttt{<=} Against a Combined \texttt{RELOP} DFA}
Suppose the DFA has transitions recognizing both \texttt{<} (alone, tagged \texttt{LT}) and
\texttt{<=} (tagged \texttt{LE}). Scanning \texttt{<=10}: after \texttt{<}, the DFA is already
in an accepting state (tagged \texttt{LT}) -- remember this position. Reading \texttt{=} still
has a valid transition (to the \texttt{LE}-tagged accepting state) -- so keep going and update
the remembered position to \emph{here} instead (a strictly longer match was just found). Reading
the next character, \texttt{1}, has \emph{no} transition from the \texttt{LE} state -- stop.
Roll back to the last remembered position (right after \texttt{=}), emit \texttt{<LE, -->}, and
resume scanning at \texttt{1}. This is exactly \emph{why} \texttt{<=} is never mis-split into
\texttt{<} followed by \texttt{=}, made completely mechanical.
\end{examplebox}

\subsection{Resolving Ties: What Happens When Two Patterns Match Equally Long}
\label{sec:l6-lexgen-ties}

\dbtag Maximal munch resolves \emph{length} ties (always prefer the longer match), but a
different kind of tie remains possible: two \emph{different} patterns matching the \emph{exact
same} lexeme, of the exact same length. The reserved-word trick from \Sref{sec:l3-beyond-db} is
one instance of this (\texttt{while} matches both a specific keyword pattern and the generic
\texttt{ID} pattern) -- and lexer generators resolve it with one simple, memorize-able rule:

\begin{coreidea}[Priority by Order of Declaration]
When several patterns match a lexeme of the \emph{same} maximal length, the lexer-generator
convention is to prefer whichever pattern was \textbf{listed first} in the original
specification. This is exactly why, in practice, keyword patterns like \texttt{if} and
\texttt{while} are always written \emph{before} the generic \texttt{ID} pattern in a lexer
specification's rule list -- listing them first guarantees the keyword-specific token wins the
tie, without needing the reserved-word symbol-table trick from \Sref{sec:l3-beyond-db} at all
(both techniques solve the same problem; a lexer generator typically uses rule ordering, while a
hand-written lexer more often uses the symbol-table trick).
\end{coreidea}

\begin{examfocus}
Rule order controls only \emph{same-length} ties. It never overrides maximal munch itself: a
\emph{longer} match from a pattern listed \emph{later} in the file still wins over a
\emph{shorter} match from a pattern listed earlier. Get the two straight: \textbf{length} beats
\textbf{position in the file}; position in the file is only the tie-breaker when lengths are
\emph{equal}.
\end{examfocus}

\subsection{Input Buffering: Reading Characters Efficiently}
\label{sec:l6-lexgen-buffering}

\dbtag Maximal munch (\Sref{sec:l6-lexgen-maximal-munch}) has one uncomfortable-sounding
consequence: the scanner sometimes reads several characters \emph{past} the eventual end of a
token before it knows to roll back (the \texttt{<=} example needed to peek at \texttt{1} before
committing to \texttt{<=}). A na\"ive implementation that requests one character at a time from
the operating system, and that can only look one step ahead, would need awkward, slow
special-casing to support this ``read ahead, then roll back'' behavior. DB's classic fix is a
purely mechanical buffering scheme that makes rollback essentially free.

\begin{defbox}{Double Buffering with Sentinels}
Maintain \emph{two} buffers of equal, fixed size (each, e.g., one disk block), laid out
contiguously in memory, plus two pointers:
\begin{itemize}
  \item \texttt{lexemeBegin} marks the start of the lexeme currently being matched;
  \item \texttt{forward} scans ahead, one character at a time, exactly as maximal munch's
  simulation loop (\Sref{sec:l6-lexgen-maximal-munch}) advances the DFA.
\end{itemize}
When \texttt{forward} reaches the end of one buffer, the \emph{other} buffer is refilled (one
system call, one full block) while scanning continues seamlessly across the boundary. Once a
token is recognized, \texttt{lexemeBegin} jumps forward to just past it, and the cycle repeats.
Rolling back \texttt{forward} to a remembered accepting position (\Sref{sec:l6-lexgen-maximal-munch})
is therefore just moving a pointer \emph{backward within already-buffered memory} -- no re-read
from disk, no special-casing at all.
\end{defbox}

A further trick avoids testing ``have I run off the end of the buffer?'' on \emph{every single}
character read (which would otherwise double the cost of the scanner's innermost loop): place a
special \textbf{sentinel} character (conventionally \texttt{eof}, distinct from every real input
character) at the very end of \emph{each} buffer. The scanner's inner loop then needs only
\emph{one} check per character -- ``did I just read the sentinel?'' -- and only in that
comparatively rare case does it do the extra work of figuring out \emph{which} buffer's sentinel
was hit (real end of file) versus a routine buffer-swap point.

\begin{examfocus}
Be able to explain \emph{why} double buffering (rather than a single buffer, or reading one
character at a time from the OS) is necessary specifically because of maximal munch's
``read-ahead-then-roll-back'' behavior -- this input-buffering material is frequently tested as
a short conceptual question (``why can't a lexer just read one character at a time from disk?'')
rather than as a diagram-drawing exercise.
\end{examfocus}

\subsection{Putting It All Together: The Generic Lexer-Generator Pipeline}
\label{sec:l6-lexgen-pipeline}

\begin{center}
\resizebox{\textwidth}{!}{%
\begin{tikzpicture}[
  box/.style={draw, thick, rounded corners, minimum height=1.05cm, minimum width=3.0cm, align=center, font=\small},
  tool/.style={box, fill=dbGreenBg, draw=dbGreen},
  data/.style={box, fill=nuLight, draw=nuNavy},
  >=Latex
]
  \node[data] (rules) at (0,0) {Rules:\\(pattern$_i$, action$_i$)\\for $i=1,\dots,k$};
  \node[tool] (combine) at (4.3,0) {Combine into\\one NFA\\(tagged accepts)};
  \node[tool] (subset) at (8.6,0) {Subset\\Construction};
  \node[tool] (min) at (12.9,0) {DFA\\Minimization};
  \node[data] (table) at (17.2,0) {Transition Table\\+ Action Table};
  \node[tool, below=1.4cm of table] (scanner) {Generated\\Scanner Code};

  \draw[->, thick] (rules) -- (combine);
  \draw[->, thick] (combine) -- (subset);
  \draw[->, thick] (subset) -- (min);
  \draw[->, thick] (min) -- (table);
  \draw[->, thick] (table) -- (scanner);
\end{tikzpicture}}
\end{center}

\begin{center}
\small\textit{Everything left of ``Generated Scanner Code'' happens once, offline, when the
lexer-generator tool runs on your rules file. What actually ships and runs on every source file
afterward is only the rightmost box: a tight loop implementing
\Sref{sec:l6-lexgen-maximal-munch}'s maximal-munch simulation over the pre-built table, driven
by the double-buffered input of \Sref{sec:l6-lexgen-buffering}, and exposed to the parser via
exactly the \texttt{getNextToken()} interface \Sref{sec:l3-where-lexical-fits} described. This
is the complete answer to ``what does a lexer-generator tool actually do'' -- Flex is one
concrete implementation of this exact pipeline, and its own \texttt{.l}-file syntax for
expressing the \emph{Rules} box is covered separately.}
\end{center}

\begin{coreidea}[Closing the Loop Back to Lecture 3's Big Picture]
Lay \Sref{sec:l3-big-picture}'s original roadmap diagram next to this section's pipeline and
they match exactly, with one addition: Lecture~\ref{lec:3} promised \emph{regular expression
$\to$ NFA $\to$ DFA $\to$ minimized DFA $\to$ transition table $\to$ lexer}, for \emph{one}
pattern. This lecture generalizes every arrow in that promise to \emph{many} patterns at once
(via the tagged-union combining step) and fills in the one detail the original diagram left
implicit -- how ``transition table $\to$ lexer'' actually happens, via maximal-munch simulation
over double-buffered input. Nothing new was invented to make this jump; the entire generalization
is the union rule from Thompson's construction, applied $k$ times instead of twice. What is
still missing -- and what Lecture~\ref{lec:7} picks up next -- is the \emph{other} way to reach
a working scanner: writing one by hand, diagram by diagram, with no generator tool at all, and
weighing that choice honestly against everything built here.
\end{coreidea}

\section{Self-Check Questions}
\label{sec:l6-selfcheck}

\begin{enrichment}[Practice -- Not Graded]
\begin{enumerate}
  \item Apply subset construction by hand to the NFA for $ab^{*}$ (built in
  Lecture~\ref{lec:5}'s self-check question 5). List each DFA state as a set of NFA states, and
  draw the resulting DFA.
  \item For the DFA you built in the previous question, apply DFA minimization. Are any states
  mergeable? Justify your answer in terms of distinguishability.
  \item Explain, in your own words, why $\varepsilon\text{-closure}$ must be applied \emph{both}
  when computing the start state \emph{and} again after every single \texttt{move}, not just
  once at the beginning.
  \item Construct a small regular expression of your own (different from
  $(a\mid b)^{*}a(a\mid b)^{n}$) where you believe the DFA will have \emph{more} states than the
  NFA, and explain your reasoning informally.
  \item Trace the maximal-munch scanning loop (\Sref{sec:l6-lexgen-maximal-munch}) on the input
  \texttt{intx} against a combined DFA for the two patterns \texttt{int} (keyword) and
  \texttt{letter(letter|digit)*} (\texttt{ID}), assuming keyword patterns are listed first. What
  single token is emitted, and why?
  \item Two patterns, \texttt{a+} and \texttt{a+b?}, are both listed in a lexer specification,
  with \texttt{a+} listed first. For the input \texttt{aaa} (with no trailing \texttt{b}
  anywhere), which pattern's action fires, and is this a length tie or a maximal-munch decision?
  \item Explain why a lexer cannot safely process input one character at a time, read directly
  from disk on every single character, given how maximal munch works -- and explain concretely
  what a sentinel character saves the scanner's inner loop from having to do on every iteration.
\end{enumerate}
\end{enrichment}

\begin{takeaways}
\begin{itemize}
  \item \textbf{Subset construction} converts any NFA into an equivalent DFA by tracking
  \emph{sets} of NFA states as single DFA states, using $\varepsilon\text{-closure}$ and
  \texttt{move} as its two building blocks; termination is guaranteed because a finite NFA has
  only finitely many state subsets.
  \item A \textbf{dead state} makes the ``no valid transition'' case explicit, and is exactly
  the signal maximal-munch scanning uses to know a match cannot be extended any further.
  \item \textbf{DFA minimization} (partition refinement) merges states that no future input
  string could ever tell apart, producing the provably smallest DFA for a language -- a
  canonical form, unique up to state renaming.
  \item NFAs and DFAs recognize \emph{exactly} the same languages, but the practical trade-off
  is real: DFA simulation is $O(1)$ per character with a pre-paid one-time build cost, while NFA
  simulation pays a per-character cost proportional to the current state set on \emph{every}
  run -- and a DFA can, in the worst case, be \textbf{exponentially larger} than its source NFA.
  \item A lexer-generator tool builds \emph{one} combined DFA out of \emph{many} token patterns
  by generalizing Thompson's union rule to $k$ patterns, tagging each pattern's own accept
  state; \textbf{maximal munch} is then just a DFA-simulation loop that remembers the last
  accepting position seen and rolls back to it; \textbf{same-length ties} between patterns are
  broken by whichever pattern was listed first.
  \item \textbf{Double buffering with sentinels} is what makes maximal munch's read-ahead
  practical: rollback becomes a cheap pointer move within already-buffered memory, and the
  sentinel collapses the scanner's per-character bounds check into a single, rare-case test.
  \item This lecture completes Lecture~\ref{lec:3}'s original big-picture promise in full:
  regular expression $\to$ NFA $\to$ DFA $\to$ minimized DFA $\to$ transition table $\to$
  working, table-driven, maximal-munch lexer -- every arrow now built, worked by hand, and tied
  back to a concrete implementation architecture. Lecture~\ref{lec:7} now looks at the same
  destination from the opposite direction: building a scanner by hand, without a generator, and
  comparing the two approaches honestly.
\end{itemize}
\end{takeaways}

\section*{References for This Lecture}
\begin{thebibliography}{99}
\bibitem[DB]{db} Aho, A.~V., Lam, M.~S., Sethi, R., and Ullman, J.~D. \textit{Compilers:
Principles, Techniques, and Tools}. 2nd~ed. Pearson, 2006.
\bibitem[ET]{et} Cooper, K.~D. and Torczon, L. \textit{Engineering a Compiler}. 2nd~ed. Morgan
Kaufmann, 2011.
\bibitem[AP]{ap} Appel, A.~W. \textit{Modern Compiler Implementation in C/Java/ML}. Cambridge
University Press, 2004.
\end{thebibliography}

\end{document}
